\section{Cost-shape generalisation: code generation on HumanEval}
\label{app:humaneval}

The following exploratory measurements use HumanEval~\cite{humaneval}
(164 Python function-completion problems). They extend the workload
coverage but expose substantial quality and protocol limitations. Pass@1
is the fraction of first completions passing the harness tests; code
validity is the fraction accepted by Python syntax parsing. Neither is
guaranteed by completing a token canvas.

\subsection{Single-pass latency and quality}

Both arms were run serial (batch-1) on the pinned stack
(\S\ref{sec:threats}), $n{=}164$ problems.

\begin{table}[ht]
\centering
\caption{Single-pass HumanEval results (serial, $n{=}164$). DG is
faster but has much lower pass@1 and syntax validity. The ms/token column
is median latency divided by mean generated tokens, not a fitted decode
slope. This table does not measure a cost curve across output lengths.}
\label{tab:humaneval}
\begin{tabular}{@{}lrrrrr@{}}
\toprule
Arm & pass@1 & code valid\% & p50 (ms) & mean tokens & ms/token \\
\midrule
DG K=8  & 0.287 & 31.1 & 280 & 99.5  & --- (staircase) \\
AR-30B  & 0.720 & 93.9 & 800 & 124.1 & 6.4 \\
AR-3B   & 0.573 & 99.4 & 560 & 107.7 & 5.2 \\
\bottomrule
\end{tabular}
\end{table}

DG's median latency is about $2.9\times$ lower than AR-30B's, but
pass@1 is 0.287 versus 0.720 and syntax validity is 31.1\% versus
93.9\%. The result therefore does not establish a quality-matched
advantage for code. In particular, the validity gap directly rules out
a general complete-canvas validity guarantee. The benchmark cannot
attribute the gap to model training data or architecture without further
controls. DiffusionGemma is a general multimodal model, not a checkpoint
trained specifically for this log-extraction study~\cite{diffusiongemma}.

\subsection{Token-budget proxy for deadlines}

The harness maps a nominal deadline $D$ to an AR token budget
$\lfloor D/\text{slope}\rfloor$. DG is evaluated with a fixed canvas
configuration. No wall-clock cancellation is enforced; observed median
latencies can exceed $D$. These are token-budget experiments, not
on-time task-success rates.

\begin{table}[ht]
\centering
\caption{Quality under nominal deadline-derived token budgets
($n{=}164$). Both models exceed 200\,ms at the median in the first
row, so DG's higher pass@1 there does not mean it meets that deadline.
A deployment comparison would count correct outputs returned before $D$
using per-event completion times.}
\label{tab:deadline}
\begin{tabular}{@{}rrrrrr@{}}
\toprule
$D$ (ms) & AR pass@1 & AR valid\% & AR p50 (ms) & DG pass@1 & DG p50 (ms) \\
\midrule
200  & 0.085 & 22.6 & 230 & 0.287 & 279 \\
400  & 0.189 & 33.5 & 434 & 0.287 & 280 \\
600  & 0.274 & 45.7 & 637 & 0.287 & 280 \\
800  & 0.311 & 54.3 & 801 & 0.287 & 280 \\
$\infty$ & 0.720 & 93.9 & 800 & 0.287 & 280 \\
\bottomrule
\end{tabular}
\end{table}

Reducing AR's token budget lowers pass@1 and syntax validity in
this harness. DG's fixed-$K$ results remain similar because the nominal
deadline does not change its decoding work. This experiment does not test
a lower-$K$ policy or establish graceful degradation under interruption.

\subsection{Correction chains: protocol caveat}

A separate generate--test--repair harness allows up to five rounds
per problem. Its recorded first-round pass fraction differs markedly
from the single-pass experiment (DG: 0.854 versus 0.287). This is a
protocol discrepancy requiring reconciliation, not evidence that the
same model suddenly improves before receiving feedback. The chain
artifacts retain outcome summaries but not the generated code needed to
fully audit those outcomes. Recorded chain latency sums model-generation
time and excludes test execution and other orchestration overhead.

\begin{table}[ht]
\centering
\caption{Recorded correction-chain summaries ($n{=}164$, up to five
rounds), pending reconciliation with the single-pass protocol. ``Cumulative''
is success by the final allowed round. First-round and total-chain
medians summarize different quantities; their ratio is not per-hop
latency growth. Chain totals include model-generation time only.}
\label{tab:chain}
\begin{tabular}{@{}lrrrr@{}}
\toprule
Arm & R1 success & cumulative & R1 p50 (ms) & total p50 (ms) \\
\midrule
DG K=8  & 0.854 & 0.957 & 287 & 305 \\
AR-30B  & 0.762 & 0.896 & 511 & 620 \\
\bottomrule
\end{tabular}
\end{table}

A total-chain median near the first-round median can occur when most
problems finish in one round, even if retries are expensive. These values
therefore do not show that DG per-hop cost remains constant or that AR
per-hop cost grows by 21\%. Nor do they validate end-to-end chain speedup
at matched quality. We retain the table for transparency but exclude it
from the principal evidence until the protocol discrepancy is resolved.
